\documentclass[12pt,a4paper]{article}
\usepackage{ctex}
\setCJKmainfont{Microsoft YaHei}
\usepackage[margin=1.5cm]{geometry}
\usepackage{tikz}
\usetikzlibrary{positioning,shapes,arrows,calc}

\begin{document}

\begin{center}
{\Large \textbf{SQL 判题考试系统 —— E-R 图}}

\vspace{1cm}

\begin{tikzpicture}[
    node distance=2.6cm and 3.0cm,
    table/.style={
      rectangle, draw=black, line width=1pt,
      minimum width=3.8cm,
      font=\small
    },
    header/.style={
      fill=black!80, text=white, font=\small\bfseries
    },
    pk/.style={font=\small\ttfamily, text=black},
    fk/.style={font=\small\ttfamily, text=black!60},
    rel/.style={font=\footnotesize, midway, fill=white, inner sep=1pt},
    line1/.style={->, >=stealth, line width=0.8pt},
    lineN/.style={->, >=stealth, line width=0.8pt},
  ]

  % ══════════════════ ROW 1 ══════════════════
  \node[table] (users) at (0, 0) {%
    \begin{tabular}{ll}
      \multicolumn{2}{c}{\textbf{users}} \\
      \hline
      \pk{user\_id}     & PK \\
      \texttt{username}   & UQ \\
      \texttt{password\_hash} & \\
      \texttt{real\_name}  & \\
      \texttt{email}      & \\
      \texttt{role}       & \\
      \texttt{status}     & \\
      \texttt{created\_at} & \\
      \texttt{updated\_at} & \\
    \end{tabular}
  };

  \node[table, right=4.8cm of users] (classes) {%
    \begin{tabular}{ll}
      \multicolumn{2}{c}{\textbf{classes}} \\
      \hline
      \pk{class\_id}    & PK \\
      \texttt{class\_name} & \\
      \fk{teacher\_id}  & FK $\rightarrow$ users \\
      \texttt{semester}  & \\
      \texttt{invite\_code} & UQ \\
      \texttt{created\_at} & \\
      \texttt{updated\_at} & \\
    \end{tabular}
  };

  \node[table, right=4.8cm of classes] (exams) {%
    \begin{tabular}{ll}
      \multicolumn{2}{c}{\textbf{exams}} \\
      \hline
      \pk{exam\_id}     & PK \\
      \texttt{exam\_name}  & \\
      \texttt{start\_time} & \\
      \texttt{end\_time}   & \\
      \texttt{duration}    & \\
      \texttt{instructions}& \\
      \texttt{lockdown}    & \\
      \texttt{status}      & \\
      \fk{creator\_id}  & FK $\rightarrow$ users \\
      \texttt{created\_at} & \\
      \texttt{updated\_at} & \\
    \end{tabular}
  };

  % ══════════════════ ROW 2 ══════════════════
  \node[table, below=1.6cm of users] (questions) {%
    \begin{tabular}{ll}
      \multicolumn{2}{c}{\textbf{questions}} \\
      \hline
      \pk{question\_id} & PK \\
      \texttt{title}      & \\
      \texttt{description} & \\
      \texttt{difficulty}  & \\
      \texttt{answer\_sql} & \\
      \fk{creator\_id}  & FK $\rightarrow$ users \\
      \texttt{schema\_sql} & \\
      \texttt{tags}        & (JSON) \\
      \texttt{visible}     & \\
      \texttt{created\_at} & \\
      \texttt{updated\_at} & \\
    \end{tabular}
  };

  \node[table, right=4.8cm of questions] (submissions) {%
    \begin{tabular}{ll}
      \multicolumn{2}{c}{\textbf{submissions}} \\
      \hline
      \pk{submission\_id} & PK \\
      \fk{user\_id}     & FK $\rightarrow$ users \\
      \fk{question\_id} & FK $\rightarrow$ questions \\
      \fk{exam\_id}     & FK $\rightarrow$ exams (NULL) \\
      \texttt{sql\_code}   & \\
      \texttt{status}      & \\
      \texttt{score}       & \\
      \texttt{runtime\_ms} & \\
      \texttt{error\_msg}  & \\
      \texttt{result}      & (JSON) \\
      \texttt{submit\_time}& \\
    \end{tabular}
  };

  \node[table, right=2.8cm of submissions, minimum width=4.2cm] (test_cases) {%
    \begin{tabular}{ll}
      \multicolumn{2}{c}{\textbf{test\_cases}} \\
      \hline
      \pk{case\_id}     & PK \\
      \fk{question\_id} & FK $\rightarrow$ questions \\
      \texttt{input\_sql}  & \\
      \texttt{expected}    & (JSON) \\
      \texttt{case\_order} & \\
      \texttt{is\_hidden}  & \\
      \texttt{created\_at} & \\
    \end{tabular}
  };

  % ══════════════════ ROW 3 (关联表) ══════════════════
  \node[table, below=0.6cm of classes, minimum width=4.2cm] (student_class) {%
    \begin{tabular}{ll}
      \multicolumn{2}{c}{\textbf{student\_class}} \\
      \hline
      \fk{student\_id} & PK, FK \\
      \fk{class\_id}   & PK, FK \\
      \texttt{status}    & \\
      \texttt{joined\_at}& \\
    \end{tabular}
  };

  \node[table, below=0.6cm of exams, minimum width=4.2cm] (exam_questions) {%
    \begin{tabular}{ll}
      \multicolumn{2}{c}{\textbf{exam\_questions}} \\
      \hline
      \fk{exam\_id}     & PK, FK \\
      \fk{question\_id} & PK, FK \\
      \texttt{score}      & \\
      \texttt{question\_order} & \\
    \end{tabular}
  };

  \node[table, below=1.2cm of exam_questions, minimum width=4.2cm] (exam_students) {%
    \begin{tabular}{ll}
      \multicolumn{2}{c}{\textbf{exam\_students}} \\
      \hline
      \fk{exam\_id}     & PK, FK \\
      \fk{student\_id}  & PK, FK \\
      \texttt{final\_score}& \\
      \texttt{status}     & \\
      \texttt{started\_at}& \\
      \texttt{submitted\_at}& \\
    \end{tabular}
  };

  % ══════════════════ 连线 ══════════════════

  % users → classes (1:N teacher creates)
  \draw[lineN] (users.east) -- ++(0.6,0) |- node[rel, pos=0.3] {1} (classes.185);
  \draw[lineN] (classes.175) |- node[rel, pos=0.7] {N} ++(-0.6,0) -- (users.east);

  % users → questions (1:N creator)
  \draw[lineN] (users.south) -- ++(0,-0.3) -| node[rel, near start] {1} (questions.10);
  % line back is implied

  % users → exams (1:N creator)
  \draw[lineN] (users.50) |- node[rel, pos=0.25] {1} (exams.170);
  % line back is implied

  % users → submissions (1:N submitter)
  \draw[lineN] (users.340) |- node[rel, pos=0.3] {1} (submissions.160);

  % classes → student_class → users (M:N students)
  \draw[lineN] (classes.south) -- (student_class.north) node[rel] {1};
  \draw[lineN] (student_class.west) -| node[rel, pos=0.4] {N} ++(-1.5,0) |- (users.300);

  % questions → test_cases (1:N)
  \draw[lineN] (questions.east) -- ++(1.2,0) |- node[rel, pos=0.25] {1} (test_cases.195);

  % questions → submissions (1:N)
  \draw[lineN] (questions.0) -- node[rel] {1} (submissions.200);

  % questions → exam_questions (M:N via exam_questions)
  \draw[lineN] (questions.15) |- node[rel, pos=0.3] {1} (exam_questions.185);

  % exams → exam_questions (1:N)
  \draw[lineN] (exams.south) -- (exam_questions.north) node[rel] {1};

  % exams → exam_students (1:N)
  \draw[lineN] (exams.355) |- node[rel, pos=0.25] {1} (exam_students.170);

  % exams → submissions (1:N, nullable exam)
  \draw[lineN, dashed] (exams.5) |- node[rel, pos=0.3] {0..N} (submissions.20);

  % users → exam_students (M:N via exam_students)
  \draw[lineN] (exam_students.west) -| node[rel, pos=0.5] {N} ++(-3.0,0) |- (users.320);

  % exam_questions to questions back
  \draw[lineN] (exam_questions.190) -| node[rel, pos=0.5] {N} (questions.5);

  % exam_students to users back
  % already drawn above

  % ══════════════════ 图例 ══════════════════
  \node[anchor=north west, font=\footnotesize] at ($(users.south west) + (0, -6.8)$) {
    \textbf{图例：}
    \quad $\rightarrow$ 1 对 N
    \quad $\dashrightarrow$ 可选关系（0 或 N）
    \quad PK = 主键
    \quad FK = 外键
    \quad 粗体下划线 = 双表联合主键
  };

\end{tikzpicture}
\end{center}

\end{document}
